\documentclass[11pt,a4paper]{article}

\usepackage[french]{babel}
\usepackage[utf8]{inputenc}
\usepackage[T1]{fontenc}
\usepackage[margin=2.5cm]{geometry}
\usepackage{graphicx}
\usepackage{booktabs}
\usepackage{hyperref}

\title{Segmentation de tissus cérébraux sur IRM \\ \large Dataset iSeg-2017}
\author{Hackathon SCIA 2026}
\date{\today}

\begin{document}

\maketitle

\section{Introduction}

Ce projet vise à segmenter des IRM cérébrales de nourrissons en trois tissus
d'intérêt (liquide céphalo-rachidien, matière grise, matière blanche) à
partir du dataset public \textit{iSeg-2017}. Ce rapport présente le dataset,
les choix de conception retenus pour préparer les données, et un problème
identifié en cours d'exploration : la présence massive de tranches
totalement vides dans le volume, ainsi que la solution mise en place pour
pouvoir les exclure et en mesurer l'impact.

\section{Dataset}

Le dataset iSeg-2017 fournit, pour chaque sujet, deux volumes IRM 3D
(modalités T1 et T2) ainsi qu'un volume de vérité terrain associant à chaque
voxel l'une de quatre classes : fond, liquide céphalo-rachidien (LCR),
matière grise et matière blanche. Les fichiers sont au format Analyze 7.5
(paires \texttt{.hdr}/\texttt{.img}).

Le dataset est scindé en deux ensembles :
\begin{itemize}
  \item \textbf{Entraînement} : 10 sujets, avec label de vérité terrain ;
  \item \textbf{Test} : 13 sujets, sans label (réservés à l'évaluation
        externe du challenge iSeg-2017).
\end{itemize}

Chaque volume 3D peut être découpé en tranches 2D selon l'un de ses trois
axes ; les tranches obtenues selon des axes différents n'ont pas
nécessairement la même taille, ce qui impose une mise à l'échelle commune
avant de pouvoir les empiler dans un même tenseur.

\section{Choix de conception}

Le projet a été structuré autour d'un pipeline en trois étapes indépendantes
du choix de segmentation proprement dit :

\begin{itemize}
  \item \textbf{Préparation} : les archives brutes du challenge sont
        extraites et réorganisées en un dossier par sujet
        (\texttt{dataset/train/subject-\emph{n}} et
        \texttt{dataset/test/subject-\emph{n}}), séparant clairement les
        sujets avec et sans label.
  \item \textbf{Chargement} : les volumes sont découpés en tranches 2D selon
        l'axe choisi, puis mis à l'échelle vers une taille carrée commune,
        soit par remplissage (padding), soit par interpolation. Le label
        étant catégoriel, son interpolation utilise systématiquement le mode
        « plus proche voisin » pour ne jamais créer de classe inexistante
        aux frontières entre tissus. Les valeurs brutes du label sont
        ensuite remappées vers des indices de classe contigus, directement
        exploitables par une fonction de coût standard de classification.
  \item \textbf{Chargement PyTorch} : les tranches sont exposées via un
        \texttt{Dataset}/\texttt{DataLoader} PyTorch classique, paramétré
        par l'axe de coupe et la modalité (T1 ou T2) utilisée en entrée.
\end{itemize}

Un module de visualisation permet d'afficher une tranche et son label
légendé par classe, ainsi que de reconstruire et afficher en 3D (via
marching cubes) les surfaces de segmentation d'un volume complet, que ce
soit à partir de la vérité terrain ou d'une prédiction de modèle.

Ces choix ne présument pas encore de l'architecture de segmentation
elle-même, qui reste à définir ; ils fixent seulement la manière dont les
données sont rendues exploitables en amont.

\section{Tranches vides : constat et solution}

En explorant le dataset, nous avons constaté qu'une grande partie des
tranches extraites sont totalement noires : ce sont les coupes prises aux
extrémités du volume 3D, en dehors de la zone occupée par le cerveau, où
l'image comme le label ne contiennent que du fond. Ces tranches n'apportent
aucune information utile à l'entraînement et risquent au contraire de le
biaiser, le modèle pouvant apprendre à prédire le fond partout pour
minimiser trivialement l'erreur.

Une mesure sur l'ensemble d'entraînement (10 sujets, découpage selon un seul
axe, soit 2560 tranches au total) confirme l'ampleur du problème :
\textbf{1550 tranches sur 2560 (environ 60~\%) sont entièrement dépourvues
de tissu segmenté}, et ce de façon homogène d'un sujet à l'autre (entre 58~\%
et 64~\% de tranches vides selon les sujets).

Pour pouvoir retirer ces tranches au besoin, et plus généralement doser la
part de tranches peu remplies conservées, nous avons ajouté au chargement
des données un seuil configurable : une tranche n'est gardée que si la
fraction de ses pixels appartenant à un tissu (donc hors fond) atteint ce
seuil. Fixé à 0, ce seuil ne filtre rien ; augmenté, il élimine d'abord les
tranches vides, puis progressivement les tranches les moins informatives.

Cette mesure suit une distribution continue plutôt qu'un simple partage
vide/non-vide : le nombre de tranches conservées chute très vite dès les
premiers seuils, avant de se stabiliser. Sur le même ensemble
d'entraînement :

\begin{center}
\begin{tabular}{lccccccc}
\toprule
Seuil minimal de tissu & 0 & 0{,}01 & 0{,}05 & 0{,}1 & 0{,}17 & 0{,}2 & 0{,}3 \\
\midrule
Tranches conservées & 2560 & 953 & 783 & 631 & 353 & 61 & 0 \\
Part du total & 100~\% & 37~\% & 31~\% & 25~\% & 14~\% & 2~\% & 0~\% \\
\bottomrule
\end{tabular}
\end{center}

Ce mécanisme de filtrage a été expérimenté dans le notebook d'exploration,
qui permet de visualiser cette distribution globalement et par sujet, et
servira de base pour comparer l'effet du retrait des tranches vides (voire
peu remplies) sur l'entraînement du modèle de segmentation.

\end{document}
